\section{Decoding}
\label{sec:decode}

There are two searches, and conflating them is the single most common source
of confusion about this system. The first turns a posteriorgram into a move
string. The second turns a move string into a move string that is
\emph{algebraically possible}. They have different objectives, different
failure modes, and --- crucially --- different standing as evidence.

\subsection{Why not peak-picking}
\label{sec:peakpick}

The obvious decoder for a per-frame onset score is: threshold, find local
maxima at least $m$ frames apart, emit one move per maximum, read the class
head at that frame. It has three failure modes, all of which were measured.

\begin{enumerate}[leftmargin=1.4em,itemsep=2pt]
\item \textbf{Merged humps.} Two turns closer than $m$ frames become one
symbol. Raising $m$ loses fast pairs; lowering it invents phantoms in the
shoulders of a single hump. In this corpus 34\% of sub-150\,ms pairs arrive
as one merged hump at the tuned operating point.
\item \textbf{Constants tuned on the wrong regime.} $m$ and the target width
$\sigma$ were both tuned on 2026-07-21 footage containing almost no fast
moves (2.5\% of pairs under 100\,ms). Against footage where 28\% of onsets
have a neighbour under 150\,ms, retuning $m: 3\to2$ and $\sigma: 2\to1$ moved
sub-150\,ms recall from 73.8\% to 83.5\%. \emph{A decoder constant tuned on
one speed regime silently becomes a ceiling in another, and from the outside
it looks exactly like a model limitation.}
\item \textbf{Premature commitment.} The candidate list is fixed before any
downstream evidence is consulted. This is the structural objection and it is
not fixable by tuning.
\end{enumerate}

\subsection{CTC prefix beam search}
\label{sec:prefixbeam}

The decoder must approximate
$\hat\ell = \arg\max_\ell p(\ell\mid x)$ where $p(\ell \mid x)$ is the sum
over alignments in \cref{eq:ctcmarginal}. Best-path decoding (per-frame
argmax, then collapse) maximises the probability of a single \emph{path},
which is a different thing, and the two differ exactly when the model is
uncertain --- which is the regime of interest. \cref{sec:main-results}
measures the gap: \GreedyGap\ points.

Prefix beam search \cite{hannun2014} maintains a set of prefixes, each
carrying two accumulators:
\begin{align}
p_b(\ell, t) &= \Pr[\text{paths over frames } 1{:}t
   \text{ that collapse to } \ell \text{ and end in blank}],\\
p_{nb}(\ell, t) &= \Pr[\text{same, ending in a non-blank symbol}].
\end{align}
The split exists to express the repeat rule. Extending prefix $\ell$ by
symbol $c$ at frame $t$:
\begin{equation}
p_{nb}(\ell c, t) \mathrel{+}= y^t_c \cdot
\begin{cases}
p_b(\ell, t-1) & \text{if } c = \mathrm{last}(\ell),\\
p_b(\ell,t-1) + p_{nb}(\ell,t-1) & \text{otherwise,}
\end{cases}
\label{eq:extend}
\end{equation}
because a repeat of the last symbol only creates a \emph{new} emission if
the previous frame was blank; otherwise it merely extends the existing
emission, which contributes to $p_{nb}(\ell,t)$ rather than to
$p_{nb}(\ell c,t)$. Staying blank carries $\ell$ forward:
$p_b(\ell,t) \mathrel{+}= y^t_\varnothing (p_b(\ell,t-1)+p_{nb}(\ell,t-1))$.
Prefixes are ranked by $\log(p_b+p_{nb})$ and pruned to the beam width (16
here). Everything is in log space: over a 1600-frame session, linear
probabilities underflow within a few hundred frames and the beam silently
degenerates to whatever ties at zero.

\paragraph{Frames are approximate.} A prefix is reached by many alignments
that disagree about which frame emitted which symbol, so each prefix keeps
the frame list of its single largest contribution. Nothing downstream uses
those frames for cost --- only for ordering and reporting --- so the
approximation never enters a score.

\paragraph{Language-model fusion, and why it is off.} Shallow fusion,
$\mathrm{rank}(\ell) = \log p_{\text{ctc}}(\ell) + \alpha \log
p_{\text{lm}}(\ell) + \beta|\ell|$, was implemented and measured. It helps
consistently across days (MER $17.3\to9.8$ and $17.1\to12.2$ by seed) and
\emph{flips sign} on the same-day validation split, because one global
$\beta$ cannot serve both regimes. It is off by default. The $\beta$ term is
not a hack: any language model systematically prefers shorter sequences, so
without an explicit per-symbol insertion bonus, fusion trades phantoms for
misses rather than reducing error.

\subsection{What CTC's independence assumption costs}
\label{sec:ctc-limits}

\Cref{eq:pathprob} factorises the path probability across frames, which
means the model has \emph{no} way to express ``having just emitted
\wca{R}, \wca{R'} is now unlikely''. All sequence-level structure has to be
carried implicitly by the trunk's receptive field --- 61 frames, about two
seconds, or roughly five turns at the corpus's median rate. Anything longer
range than that is invisible to the objective.

For this domain that is a real limitation, because the output has enormous
long-range structure: the last third of a solve is drawn from a repertoire
of a few dozen memorised words, and a solver's move distribution is far from
i.i.d. Two remedies exist and both were tried.

\begin{itemize}[leftmargin=1.3em,itemsep=2pt]
\item \emph{Shallow fusion with a move language model} (above). Helps
cross-day, hurts same-day, sign-unstable --- off by default.
\item \emph{An explicit algorithm prior} in the group decoder --- recognise a
known last-layer word and insert it as a single beam transition. Measured
inert in the forward direction and net negative in the backward direction
(\cref{sec:negatives}).
\end{itemize}

The honest reading is that neither failed because sequence structure is
absent --- it is obviously present --- but because both were applied
\emph{after} the frame-level model had already lost the evidence. A model
that conditions on its own emission history (an RNN-transducer or an
attention decoder) would carry the prior where it can still change what is
recognised. That is the most defensible architectural change available and
it has not been tried.

\subsection{The group-theoretic decode}
\label{sec:reconstruct}

\paragraph{Setup.} We have a start state $s_0 \in G$ (the server-issued
scramble, known exactly), an end state (solved, $e$), and a detected
sequence of $n$ onsets, each with a 12-way posterior $p_j$ over classes.
We want the most probable move word $w$ such that $s_0 \cdot w = e$.

Model the recogniser as a noisy channel with three error types --- the
alignment operations of \cref{sec:experiments}: a \emph{substitution} (onset
detected, named wrong), a \emph{phantom} (spurious onset), a \emph{miss}
(onset never emitted). A hypothesis is a sequence of edit decisions; its cost
is a sum of negative log-probabilities, so minimising cost is MAP estimation.

\paragraph{The cost model.} For onset $j$, accepting class $k$ costs
\begin{equation}
c_j(k) \;=\; \log \max_i q_j(i) \;-\; \log q_j(k),
\qquad
q_j \;=\; (1-\rho_j)\,p_j \;+\; \tfrac{\lambda_u}{12}\mathbf{1}
        \;+\; \lambda_i u_j\, \delta_{\mathrm{inv}(\arg\max p_j)},
\label{eq:acceptcost}
\end{equation}
where $u_j = 1 - \max_i p_j(i)$ is the onset's uncertainty and
$\rho_j = \lambda_i u_j + \lambda_u$. Accepting the model's own argmax is
free by construction; everything else is priced by how far the smoothed
posterior falls short of it. Two smoothing terms are added because the raw
softmax is badly calibrated in a specific, measured way: it is overconfident
precisely on the temporal inverse (\wca{U} versus \wca{U'}), so
$\lambda_i = 0.20$ of prior mass, scaled by the onset's own uncertainty, is
moved onto the inverse class, plus a flat $\lambda_u = 0.04$ uniform floor.

Deleting onset $j$ (declaring it a phantom) costs
$C_{\mathrm{del}} + w\cdot(\text{onset strength})$ with
$C_{\mathrm{del}}=2.6$. Making deletion depend on strength is not a
refinement, it is load-bearing: with a flat deletion cost the beam drowns in
``delete something here, insert something there'' pairs --- there are
$O(n^2\cdot 12)$ of them, they are parity-balanced, and they are all cheaper
than a true multi-substitution story. Onset strength is genuine evidence
against deletion and pricing it removes the flood.

Inserting a move (recovering a miss) costs a flat $C_{\mathrm{ins}} = 4.0$,
roughly $-\log$ of the miss odds plus the 12-way choice of which move.

\paragraph{Why beam search and not exact search.} At the measured error rate
an exact IDA* over edit cost is hopeless, because the constraint that kills
wrong hypotheses --- ``does this reach solved'' --- only pays off at the
\emph{end} of the sequence. At a budget sufficient for ten edits the tree
explodes long before pruning bites. So: keep the $K$ cheapest hypotheses per
onset ($K = 4000$, retried at $16000$), deduplicated by cube state, ranked
by cost-so-far plus an admissible lookahead.

\paragraph{The lookahead.}
\label{sec:heuristic}
Three components, two of them true lower bounds.
\begin{itemize}[leftmargin=1.3em,itemsep=2pt]
\item \emph{Pattern databases.} Breadth-first tables over the corner-twist,
edge-flip and UD-slice coordinates give an admissible quarter-turn distance
from any state to solved; the max over tables is used.
\item \emph{Parity.} Every quarter turn flips permutation parity
(\cref{sec:background}), so a hypothesis whose achievable remaining move
count has the wrong parity provably needs a further insertion or deletion.
\item \emph{A suffix-consistency ladder.} This one is the interesting part.
Precompute $S_j$, the group product of the \emph{remaining} detected moves as
classified. A hypothesis is consistent with the rest of the detections iff
$\mathrm{rel} = \mathrm{state}\cdot S_j$ is the identity. Grading ``almost
consistent'' by pattern distance does \emph{not} work --- a conjugated error
scrambles the coordinate metrics, so $h(\mathrm{rel})$ measures $20\pm3$ for
\emph{any} hypothesis with at least one unfixed error regardless of how many,
and feeding that noise into the ranking evicts the true hypothesis. What
works is exact algebra: $\mathrm{rel}$ is repairable by one future insertion
of $m$ at gap $q$ iff $\mathrm{rel} = S_q^{-1} m^{-1} S_q$, by one future
deletion of detection $q$ iff $\mathrm{rel} = S_{q+1}^{-1} d_q S_{q+1}$, and
by one future substitution of $m$ for $d_q$ iff
$\mathrm{rel} = S_{q+1}^{-1}(m^{-1}d_q)S_{q+1}$. Those $\sim\!25n$ elements
go into a hash set; the ranking adds a weight per ``level'' (0 = consistent
now, 1 = one edit away, capped otherwise). It is noise-free, and the moment a
hypothesis has repaired everything but one discrepancy it out-ranks the
cheap procrastinators an absolute-cost beam would bury it under. It is a
ranking term, never a hard prune.
\end{itemize}

\paragraph{The tail is exact.} After the last onset, each candidate end state
is finished by iterative-deepening search over words of length
$\le M_{\mathrm{end}} = 4$, pruned by the full pattern-database bound. This
covers end-of-solve moves the camera never saw. \textbf{The small budget is
what keeps the whole thing sound}: every state is within 26 quarter turns of
solved, so with a generous budget \emph{every} branch ``verifies'' and the
result means nothing.

\subsection{What the group decode can and cannot repair}
\label{sec:decode-limits}

The endpoint constraint is \emph{global and terminal}. It discriminates
nothing until the last onset. A session that cannot complete a consistent
path to solved therefore receives only generic smoothing, and all of the
decode's power concentrates in sessions that were nearly right already. This
is why the measured contribution is about a point (\cref{sec:decode-results})
rather than the large number the framing suggests.

Per channel:
\begin{itemize}[leftmargin=1.3em,itemsep=2pt]
\item \textbf{Substitutions} are the channel the constraint fixes nearly for
free --- the group pins down which move a mis-named onset must have been. It
is also the \emph{smallest} channel (\SubShare\% of error mass).
\item \textbf{Phantoms} are cheap to delete when the onset was weak.
\item \textbf{Misses} force the decoder to invent a move with no posterior
evidence at all: a 12-way guess whose only support is an endpoint dozens of
moves away. They are \MissShare\% of the error mass, i.e.\ most of it.
\end{itemize}

The envelope is best stated in \emph{cost}, not error count, because cost is
what actually binds. Measured on synthetic corruptions and on a real replay:
a true story costing ${\sim}6$ survives 90 onsets at beam 4000; one costing
${\sim}10$ does not, and each insertion is a flat 4.0. So there is roughly
\emph{one insertion of headroom beyond whatever the recogniser's own mistakes
already cost.} Buying down recognition error buys miss recovery for free;
widening the beam does not.

That envelope is the whole story of \cref{sec:decode-results}: on the
held-out set the cheapest true story costs \MinGtPathCost\ and the median
costs \MedGtPathCost, so nothing is within reach and the decode's measured
contribution is \DecodeGain\ points. A decoder with exactly one anchor, at
the very end, cannot help a sequence that needs a dozen insertions to reach
it.
